% Part: model-theory
% Chapter: interpolation
% Section: introduction

\documentclass[../../../include/open-logic-section]{subfiles}

\begin{document}

\olfileid{mod}{int}{int}

\olsection{سريزه}

د منځګړيتوب قضيه لاندې پايله ده: فرض کړئ
$\Entails !A \lif !B$۔ بيا داسې !!a{sentence}~$!C$ شته چې
$\Entails !A \lif !C$ او $\Entails !C \lif !B$ وي۔ سربېره پر دې،
په~$!C$ کښې هر !!{constant}، !!{function} او !!{predicate} (له
$\eq$ پرته) په~$!A$ او~$!B$ دواړو کښې راځي۔ !!{sentence}~$!C$ د
$!A$ او~$!B$ يو \emph{منځګړی} بلل کېږي۔

د منځګړيتوب قضيه په خپل ذات کښې په زړه پورې ده، خو اصلي ارزښت ئې
په دې کښې دے چې د يوې تيورۍ دننه د تعريف‌وړتيا په هکله پايلې، او
هغه شرطونه پرې ثابتولے شو چې لاندې ئې د دوو سازګارو تيوريو يوځاے
کول يوه سازګاره تيوري ورکوي۔ لومړۍ پايله د بېت د تعريف‌وړتيا قضيه،
او دويمه د رابنسن د ګډې سازګارۍ قضيه بلل کېږي۔

\end{document}
